% Figure for Oscillations, Waves and Optics §A7.3 (N-slit pattern normalised to its principal maxima, I/I_max = sin^2(N phi/2)/(N^2 sin^2(phi/2)), computed for N = 2 (thin blue), N = 5 (green) and N = 20 (thick red): as N grows the principal maxima at d sin(theta) = m lambda stay put but become narrow, and the secondary maxima fade).
\begin{tikzpicture}[font=\small]
  \begin{axis}[nfig, clip=false, width=11.5cm, height=5.2cm, xmin=-1.35, xmax=1.75, ymin=0, ymax=1.18,
      xlabel={$d\sin\theta/\lambda$}, ylabel={$I/I_{\max}$},
      xtick={-1,-0.5,0,0.5,1}, xticklabels={$-1$,$-\tfrac12$,$0$,$\tfrac12$,$1$}, ytick={0,0.5,1},
      xlabel style={at={(axis description cs:1.0,0)}, anchor=west},
      ylabel style={rotate=-90, at={(axis description cs:0,1)}, anchor=south}]
    \addplot[figB, thin, domain=-1.3:1.3, samples=600] {cos(deg(pi*x))^2};
    \addplot[figG, domain=-1.3:1.3, samples=1200] {(sin(deg(5*pi*(x+0.000137)))/(5*sin(deg(pi*(x+0.000137)))))^2};
    \addplot[figR, line width=1.3pt, domain=-1.3:1.3, samples=2400] {(sin(deg(20*pi*(x+0.000137)))/(20*sin(deg(pi*(x+0.000137)))))^2};
    \node[font=\scriptsize, figB, anchor=west] at (axis cs:1.35,0.85) {$N = 2$};
    \node[font=\scriptsize, figG, anchor=west] at (axis cs:1.35,0.55) {$N = 5$};
    \node[font=\scriptsize, figR, anchor=west] at (axis cs:1.35,0.25) {$N = 20$};
  \end{axis}
\end{tikzpicture}
